\chapter{Literature Review}

% Stored products refer to raw and processed non-perishable commodities, including grains, legumes, nuts, spices, and manufactured food. Fundamental for food security and nutrition, they are also economically valuable and represent a large portion of global trade. These products are susceptible to infestation by various pests, resulting in quantitative and qualitative losses as well as pose a risk to human health and safety. Efficaciously inspecting products for pests throughout the global supply chain is not only crucial for the integrity of the product and mitigate greater infestation but also secure regional agricultural and environmental ecologies from the introduction of invasive species.

\section{Stored Product and Insect Pests}

Grains include cereals, pulses, and oilseeds. Cereals are grasses cultivated for its edible grain such as wheat, rice, corn, rye, millet, and barley \cite{awikaMajorCerealGrains2011}. Pulses are the edible seeds of the \textit{Leguminoseae} family and include beans, lentils, peas, and chickpeas. Peanuts, soybeans, and alfalfa are legumes but are considered to be oilseeds, such as sunflower seeds, rapeseeds, and sesame seeds \cite{sathiyaseelanMajorInsectPests2022}. Stored products also include nuts, spices, sugar, dried fruits, and processed food \cite{neethirajanDetectionTechniquesStoredProduct2007}.

Insects attack stored products in a variety of ways and are categorized by their feeding behavior. The rice weevil, \textit{Sitophilus oryzae}, the granary weevil, \textit{Sitophilus granarius}, and the lesser grain borer, \textit{Rhyzopertha dominica}, are examples of internal feeders, which bore into the kernels of grains and seeds as larvae and feed on the endosperm, visually undetectable, until they emerge as adults leaving burrowing holes and powder residue. External feeders, such as the drugstore beetle, \textit{Stegobium paniceum}, the cigarette beetle, \textit{Lasioderma serricorne}, and the khapra beetle, \textit{Trogoderma granarium}, chew through the outer seed coat of the grain exterior and then devour the inside. The red flour beetle, \textit{Tribolium castaneum}, and the sawtoothed grain beetle, \textit{Oryzaephilus surinamensis} are examples of scavenger feeders, which prefer broken down and processed material such as split kernels and flour. Secondary pests, like the mealworm, \textit{Tenebrio molitor}, eat deteriorating, damp, or moldy material \cite{bennettTrumansScientificGuide1988}.

Stored product pests (SPPs) can infest crops at any point from the field to storage. Infestations lead to both quantitative and qualitative losses. Direct consumption of grain reduces its weight and volume, while contamination from excrement, webbing, and carcasses degrades its quality. By consuming grain embryos, insects lower the protein content and alter the chemical composition of the kernels. This damage promotes infections, attracts other pathogens, and reduces seed viability. Furthermore, SPPs can damage storage containers and equipment \cite{devipriyar.InsectPestDetection2022}. The rate of infestation is accelerated by high temperatures, and the moisture and heat generated by the insects create an even more favorable environment for their proliferation \cite{muller-blenkleMethodAcousticStorage2023}.

An unmitigated infestation can lead to significant ecological, agricultural, and economic impact. Estimates by the Food and Agriculture Organization (FAO) of the United Nations suggest that \SI{20}{\percent} of the global food production is lost to herbivorous insects \cite{mohamedn.sallamInsectDamagePostHarvest2013}. Other sources state that post-harvest insects are responsible for \SIrange{10}{20}{\percent} of storage losses \cite{phillipsBiorationalApproachesManaging2010,gogiToxicityRepellencePlant2016} and cost up to \$470 billion USD per year in 2001 \cite{aliSustainableManagementInsect2023}. Introduction of non-native insects into ecological environments can further exacerbate negative effects if the species has no natural predators. Invasive species cost the United States \$2 trillion USD from 1960 to 2020, costing \$1 billion USD a year from 2010 to 2020 in damages, losses, and prevention \cite{fantle-lepczykEconomicCostsBiological2022}. Globally, stored product pests inflict over \$100 billion USD \cite{harmanIncreasingGlobalRisk2025}.

The khapra beetle, \textit{T. granarium}, also known as the cabinet beetle, represents one of the most significant stored product threats. It is capable of destroying \SIrange{30}{70}{\percent} of an infested commodity. Its larvae can survive years without food through diapause, are highly resistant to insecticides, prefer hiding in cracks and crevices of storage containers, and produce an abundance of exuviae and frass during their infestation. Native to South Asia, it has not become established in the United States. The United States Department of Agriculture (USDA) estimates that an established infestation would result in economic consequences up to billions of dollars from direct crop losses and immediate embargoes from trading partners. In the 1950s, the US spent 13 years eradicating infestations at 600 sites across southwestern states with a cost of \$141.8 million USD \cite{harmanIncreasingGlobalRisk2025}. As a result, the US maintains a zero-tolerance policy and enforces strict phytosanitary restrictions and inspections for products that might contain the khapra beetle \cite{u.s.departmentofagricultureanimalandplanthealthinspectionserviceKhapraBeetleProgram2022}.

\newpage
\section{Inspection Methods for Stored Product Insect Infestation}

\subsection{Standards for Inspecting Stored Product for Pests}

The International Organization of Standards (ISO) standard 6639 guides the means of storing cereals and pulses as well as controlling and detecting attacks by pests. It specifies that a container is considered infested if an insect is found within \SI{1}{\kilo\gram} of the material sample. During inspection, the collected material sample is divided into \SIrange{200}{300}{\gram} test samples. If insects are not found, the material is considered to be free of infestation \cite{ISO6639Cereals1986}. The USDA also guides the Customs and Border Protection Agriculture Specialists (CBPAS) to deem a container infested if one or two insects are found within \SI{1}{\kilo\gram} of material sample \cite{service2015a}.

\subsection{Methods for Inspecting Stored Products for Insects}

The main method employed at United States Ports of Entry to inspect stored product for insects is through manual inspection, where the chosen commodity is emptied on an inspection table and personnel, often with tools including sieves, tweezers, magnifying glasses, and intense lights, sift through the material looking for live and dead insects at different life stages, fragments, or signs of insect infestation such as webbing, casings, frass, or burrowing holes in the kernels. This process is time and physically consuming, requiring experienced personnel or with a background in entomology to identify the species of the insect in order to determine the appropriate course of action. If the insect is identified as a quarantine pest, the material is either treated, destroyed, or rejected and a Emergency Action Notification (EAN) is issued to the importer \cite{service2015a}.


The ISO 6639 standard describes several methods for detecting insects in stored grain products, focusing on carbon dioxide analysis, the ninhydrin method, the whole grain flotation method, as well as acoustic and X-ray methods \cite{ISO6639Cereals1986}. Various tools and devices have been developed, categorized by the sensing technologies they leverage, in order to expedite areas within the inspection including automating the examination process, detecting insect infestation, or identifying the insect species. The application of these inspection methods has to balance the accuracy of detecting and classifying pests of various life stages with safeguarding product quality, minimizing associated costs, ensuring the efficiency of supply chain movement, and limiting product damage and loss during examination.

\subsubsection{Physical Methods}

The earliest form of stored product insect detection is through physical detection, specifically the Berlese-Tullgren trap. The biologist Antonio Berlese presented his eponymous funnel in 1905. In 1905, Albert Tullgren modified the design by introducing an electric light bulb as the heat source and increasing the gradient through an iron sheet drum.

Probe traps can be placed at different grain depths for extensive amounts of time and are periodically removed to determine the number, species, and density of the infestation. This method is restricted to easily accessible locations within the container, limits the temporal availability of the results, and the capture rate is influenced by the insect species, the grain type, as well as environmental conditions including temperature and humidity \cite{neethirajanDetectionTechniquesStoredProduct2007}.

The Electronic Grain Probe Insect Counter (EGPIC) automates the process of counting insects captured by a probe trap inserted into the grain. The probe consists of a series of electrodes that create an electric field. As insects pass through the probe, they disrupt the field, generating electrical signals that are detected and counted by the system \cite{litzkowElectronicGrainProbe1997, shumanCommercializationElectronicGrain2001}.

Rolling sieves, such as the Insectomat 5K, or electrical conductance machines speed up and automate the detection process but are expensive and perform poorly with infested kernels, insect eggs, and young larvae.


Physical methods use probe traps and counters, such as the Berlese-Tullgren trap, the TNAU Insect Detection System, and the Electronic Grain Probe Insect Counter (EGPIC), to enumerate the captured insects and determine the level of infestation using extrapolation. Some probe traps can be placed at different grain depths for extensive amounts of time and periodically removed to evaluate the number, species, and density of the infestation. Although low-cost, passive, and effective at long-term storage facilities, this method is restricted to easily accessible locations within the container and the capture rate is influenced by the insect species, grain type, as well as temperature and humidity. Although visual lures and pheromones assist in attracting insects to the traps, these physical methods require a long time duration and cannot detect insects buried inside kernels \cite{neethirajanDetectionTechniquesStoredProduct2007, hoffmanStoredProductPestMonitoring2005}. Other physical methods include rolling machines that automate the sieving of products, systems that monitor for changes in electrical conductance between whole, damaged, and infested kernels, and the flotation method which separates possible insects from product by density, but these tend to damage the material in the process \cite{bonjourEvaluatingRemoteMonitoring2004, t.pearsonDetectionWheatKernels2007, thindDeterminationLowLevels2000}.

\subsubsection{Chemical Methods}

Chemical sensors monitor the $\text{O}_2$ and $\text{CO}_2$ levels in the stored product to detect the respiration of the insects or by measuring for uric acid, the main element of insect excreta \cite{rajeshEfficacyTnauStackProbe2015, srivastavaFuzzyControllerBased2019, srivastavaProbabilisticArtificialNeural2019}. Even though it was not developed for insect detection, the electronic nose (E-nose), which consists of gas sensors with varying sensitivities and pattern recognition software, was able to detect the volatile compounds and react to the changing electrical properties to identify an infestation of the lesser grain borer, \textit{R. dominica} in wheat \cite{srivastavaFuzzyControllerBased2019,srivastavaProbabilisticArtificialNeural2019}.

\subsubsection{Spectral Methods}

Advanced imaging technologies offer non-destructive solutions for pest detection. By using an X-ray source and imaging system, inspectors can visualize insects concealed within stored grains, fruits, and vegetables, as well as find damaged products \cite{XRayInspectionSystems2023, XRayDetectionFood2023}. Alternatively, near-infrared (NIR) spectroscopy leverages the unique light absorption properties of materials to detect insects and classify infested kernels, though its need for constant recalibration is a drawback \cite{neethirajanDetectionTechniquesStoredProduct2007, elizabethb.maghirangDetectingSingleWheat2002, keszthelyiNonDestructiveImagingSpectroscopic2020}. Thermal imaging detects the infrared energy emitted by objects, creating heat maps that expose infestations through temperature anomalies \cite{bangaTechniquesInsectDetection2018}. To boost precision, these systems are often paired with deep learning algorithms \cite{devipriyar.InsectPestDetection2022}.

\citeauthor{mankinMicrowaveRadarDetection2004} demonstrated that a \SI{24}{\giga\hertz} microwave Doppler radar, originally used for termite detection, can also find stored product pests. The insect movement is detected by radiating a microwave signal and observing the change in the frequency of the received reflection, which increases as the insects move slightly closer or decreases as they move further away from the sensor. The system was able to detect the cigarette beetle, \textit{Lasioderma serricorne}, the sawtoothed grain beetle, \textit{Oryzaephilus surinamensis}, the black carpet beetle, \textit{Attagenus unicolor}, and the red flour beetle, \textit{Tribolium castaneum} within a corn meal and flour mixture \cite{mankinMicrowaveRadarDetection2004}.

% Acoustic sensors leverage traditional microphones, piezoelectric discs, and accelerometers to detect the sound waves and vibrations generated by the insects as they move, communicate, masticate, and procreate in the material \cite{mankinAcousticIndicatorsTargeted2010, mankinAutomatedApplicationsAcoustics2021,flynnAcousticMethodsInvasive2016,kadyrovVibroAcousticSignaturesVarious2024}. These systems are designed as standalone receptacles where product is placed momentarily during examination or as  probes and sensors that could be inserted into the product or mounted on the storage containers even for continuous monitoring. Microwave systems track the movement of insects through a material by detecting differences between emitted and reflected radar frequencies \cite{mankinMicrowaveRadarDetection2004}. Although both these methods are relatively low-cost and non-destructive, they are highly susceptible to external environmental noise and vibration.

\newpage
\section{Acoustic Methods for Detecting Insects in Stored Products}

\subsection{Acoustic Characteristics}

Insects generate sounds within the material during locomotion through friction between the body and the material as well as through the movement of the wings, legs, and other body parts. Sounds are also generated by mastication, communication, and procreation. The production and magnitude of these sounds are strongly affected by the insect size and lifestage, feeding habits, speed, activity, the density of the infestation, and the general behavior of the subject which can vary with temperature and external disturbance. The propagation of the sound is dependent on the material properties including density, kernel size and hardness, relative humidity, and the amount of air in the material \cite{mankinPerspectivePromiseCentury2011}. The distance between the insect and the sensor as well as variation in the recording equipment and environmental conditions can also introduce differences in the results from different studies and publications. Often the sensors, signal features, processing, and algorithms are common with the platforms, analysis, and detection methods for insect infestations in other materials such as plants, soil, and wood \cite{sutinAutomatedAcousticDetection2019}.

The insect activity manifests in acoustic recordings as impulses, high magnitude signals across a wide frequency range and short duration which can be characterized by their intensity, duration, and spectral features. Microphones, piezoelectric sensors, laser vibrometers, micromechanical sensors (MEMS), accelerometers, ultrasonic transducers, and seismic sensors have all been effective at capturing these features to be used by detection algorithms to determine the presence of an insect, monitor for hidden infestations without destroying the product, classify specific species and their respective life stage, quantify the magnitude of the infestation, as well as mitigate the influence of external noise.

The spectral domain refers to the frequency content of a sound. Generally, insect noise occurs in the frequency band between \SIrange{80}{8000}{\hertz} \cite{mankinPerspectivePromiseCentury2011} and can vary between the species and lifestages \cite{azizifarsaniBioacousticsTrogodermaGranarium2023,fleurat-lessardAcousticDetectionAutomatic2006}. The khapra beetle, \textit{T. granarium}, in wheat grains, was most dominant in frequencies \SIrange{1000}{8000}{\kilo\hertz}. The adult and larva occupied different bands, \SIrange{3500}{4500}{\hertz} and \SIrange{1000}{1700}{\hertz}, respectively. Female adults dominated mainly in the \SIrange{1000}{3000}{\hertz} and \SIrange{4000}{5000}{\hertz} frequency bands \cite{azizifarsaniBioacousticsTrogodermaGranarium2023}. The peak energies for an adult and larva of the rice weevil, \textit{Sitophilus oryzae}, in wheat grains were within \SIrange{1800}{3000}{\hertz} and \SIrange{1300}{2000}{\hertz}, respectively \cite{fleurat-lessardAcousticDetectionAutomatic2006}. This variety in frequency content enabled \citeauthor{schwabAutomaticAcousticalSurveillance2005} to distinguish between the larval and adult stages of the \textit{S. granarius}, \textit{T. confusum}, and the \textit{R. dominica} through spectral features alone.

The temporal domain refers to the time-varying nature of the signal including amplitude and duration. A single impulse might be \SIrange{3}{30}{\milli\second}. Bursts, which are a series of pulses that can be enumerated by their number and rate within a time period, can range from 50 to 2000 impulses for a duration from \SIrange{3}{30}{\milli\second} with up to \SI{250}{\milli\second} intervals of separation. Adult lesser grain borers, \textit{R. dominica}, and red flour beetles, \textit{T. castaneum}, produced 37 and 80 times more sound than their larvae \cite{mankinPerspectivePromiseCentury2011}. An increase in impulses could be indicative of a higher population infesting the material \cite{hagstrumAutomatedAcousticalMonitoring1991}. Some algorithms count the number of bursts and impulses to determine an infestation \cite{mankinAcousticIndicatorsTargeted2010,mankinPerspectivePromiseCentury2011, njorogeAcousticPitfallTrap2019, mankinAutomatedApplicationsAcoustics2021}.

Signal processing such as filters and transformations can be used to isolate the frequency range of interest and expose the relevant signals while removing superfluous noise from analysis that can vary between sensors and environments. The waveform envelope allows for threshold crossing algorithms to determine an insect presence if the amplitude exceeds a statistically or empirically determined cutoff \cite{mankinAutomatedApplicationsAcoustics2021}. \citeauthor{mankinNoiseShieldingAcoustic1996} observed that the sounds of \textit{S. oryzae} movement and feeding were \SI{23}{\deci\bel} \cite{mankinNoiseShieldingAcoustic1996}. \citeauthor{balingbingDeterminingSoundSignatures2024} continuously recorded the \textit{R. dominica}, the \textit{S. oryzae}, and the \textit{T. castaneum} in rice over a \SI{24}{\hour} time frame using an Adafruit 12S MEMS SPH0645 microphone and calculated the root-mean-square (RMS) of \SI{3}{\minute} windows at each hour of a signal filtered between \SIrange{93.75}{2500}{\hertz} \cite{balingbingDeterminingSoundSignatures2024}. \citeauthor{eliopoulosDetectionAdultBeetles2015} utilized a \SI{14}{\centi\meter} long piezoelectric sensor inserted with a probe into hard wheat grain. The system was able accurately detected the presence of an insect with a probability range from \SIrange{72}{100}{\percent}, a \SI{100}{\percent} false positive rate. Their algorithm was based on extracting the Hilbert envelope and indicating a detection when it meant a certain threshold \cite{eliopoulosDetectionAdultBeetles2015}. \todo{Something is wrong with that 100 false positive rate}

The regularity and duration of the monitoring also has an effect on the performance of the sensor and its algorithm. \citeauthor{vickSoundDetectionStoredProduct1988} noted that the \textit{R. dominica}, \textit{S. cerealella}, and \textit{S. oryzae} produced sounds \SI{61}{\percent}, \SI{71}{\percent}, and \SI{90}{\percent} of the time in grain during sampling durations of \SI{180}{\second} \cite{vickSoundDetectionStoredProduct1988}. Short analysis durations can increase high rates of false negatives especially if the insects are not active during the sampling period. Longer sampling periods, such as \SIrange{10}{20}{\minute} allow for a more accurate detection as well as enable identification of background noise. Deployment of multiple sensors simultaneously could increase detection performance, determine locations of insects, level of infestation, and to offset the frequency and duration of monitoring \cite{mankinPerspectivePromiseCentury2011}. \citeauthor{hagstrumAutomatedAcousticalMonitoring1991} concluded that the frequency of monitoring a single sensor had an equivalent performance as monitoring multiple sensors at a lower frequency \cite{hagstrumAutomatedAcousticalMonitoring1991}.

Limiting the influence of external noise and vibration from the surrounding environment, such as factory machinery, vehicles, and other sources,  is critical to the success of these algorithms. External noise can mask the low-intensity sounds produced by insect feeding and movement, which typically range from \SIrange{20}{80}{\deci\bel} SPL (sound pressure level) at frequencies below \SI{10}{\kilo\hertz}. Often this background noise can be reduced physically through shielding, sound-proofing, and vibration isolation. \citeauthor{mankinNoiseShieldingAcoustic1996} was able to detect internally feeding larvae in grain samples at an inspection facility with high noise backgrounds by shielding the acoustic detectors using a \SI{81}{\kilo\gram}, \SI{26}{\centi\meter} thick, six layer muffle box that attenuated sound by \SIrange{70}{85}{\deci\bel} between \SIrange{1}{10}{\kilo\hertz} \cite{mankinNoiseShieldingAcoustic1996}. Other methods could be through box in a box design, anechoic chambers, employing Faraday cages, and shock mounts. Alternatively, reference sensors could be used to measure and identify peripheral noise and vibration \cite{mankinAcousticIndicatorsTargeted2010}.

The spectral and temporal differences between the insects and background noise can also be used to distinguish between the two. \citeauthor{mankinAcousticCharacteristicsDynastid2009} used filtering to reduce signals below \SI{200}{\hertz}. In another algorithm, the bursts characteristic to distinguish the insect sounds from the continuous or isolated impulses commonly found in background noise \cite{mankinDetectionAnoplophoraGlabripennis2008, mankinAcousticCharacteristicsDynastid2009}.

Machine learning algorithms can analyze acoustic signals and identify patterns associated with insect activity as well as distinguish external noise however they depend on large amounts of training data. The ML models can be trained on the raw audio signals or on extracted features and the data can be provided labelled or unlabelled \cite{kohlbergBuzzesBytesSystematic2024}. These algorithms include various nueral network classifiers such as convolution neural networks (CNN), probabilistic neural networks (PNN), perceptual learning prediction (PLP), decision trees and forests, hidden Markov models (HMM), support vector machines (SVM), and bayesian classifiers \cite{mankinAutomatedApplicationsAcoustics2021}. Deep learning combines multiple layers of neural networks to learn the features and patterns of the data. Since machine learning algorithms are data-driven, the performance of the model might not generalize to new species or materials, if the training data is not diverse enough.

\citeauthor{phungAutomatedInsectDetection2017} used the BugBytes dataset to compare training a linear discriminant, fine kNN, linear SVN, and Bagged Tree model using low-level features such as the zero-crossing rate, short-time average energy, bandwidth, and cut-off frequency as well as the mel-frequency ceptstrum coefficients (MFCC). They concluded that the Bagged Tree model provided the best accuracy for species classification and insect classification with an accuracy of \SI{97.1}{\percent} and \SI{92.3}{\percent}, respectively while using both low-level and MFCC features. \citeauthor{montemayorDetectingRiceWeevils2024} used the mel-frequency cepstrum coefficients from the audio signals and their log-scaled representations as Time-Frequency Patches (TF-Patches) as inputs to a Convolutional Neural Network (CNN) \cite{montemayorDetectingRiceWeevils2024}. \citeauthor{faissAdaptiveRepresentationsSound2023} observed that training a LEAF model with raw waveform data from the diverse InsectSet32, InsectSet47, and InsectSet66 resulted in better performance than a CNN model with spectrograms, a graphical visualization of the frequency magnitude over time \cite{faissAdaptiveRepresentationsSound2023}.

\subsection{Acoustic Detection Systems}

The Acoustic Location Fixing Insect Detector (ALFID) system was designed to be vertically mounted with grains gravity loaded and unloaded. The system utilizes an array of sixteen linearly oriented piezoelectric sensors ranged to identify the number of feeding larvae and their location within the sample. After amplifying the signals from the sensors by \SI{80}{\deci\bel}, a bandpass filter from \SIrange{1000}{10,000}{\hertz} is applied. The arrival time of the signal is identified using amplitude threshold detection and the location is determined by a detection order algorithm.The system was able to detect the presence of up to three rice weevil, \textit{S. oryzae}, larvae in \SI{1}{\kilo\gram} of wheat with a detection probability of \SI{64}{\percent} and a false positive rate of \SI{8}{\percent} \cite{shumanQuantitativeAcousticalDetection1993,vickAcousticLocationFixing1996}.

The Beetle Sound Tube was developed to support semi-permanent monitoring of silos and flat stores of grain to detect low levels of infestation early and provide remote reports. The system comprises of perforated metal tubes that are installed into the grain mass. Each tube features an insect trap that can be inspected for species classification, an acoustic sensor, and a climate sensor to measure changes in temperature and humidity. The signal processing on the acoustic signals utilize $\frac{1}{3}$-octave analysis, calculating the energy level for 16 $\frac{1}{3}$-octave bands from \SIrange{250}{8000}{\hertz}. The algorithm determined an insect was present if the energy level exceeded a threshold of \SI{25}{\deci\bel} within \SIrange{250}{800}{\hertz}, \SI{20}{\deci\bel} within \SIrange{800}{2500}{\hertz}, or \SI{15}{\deci\bel} within \SIrange{2500}{8000}{\hertz}, above the background noise. The sensor and algorithm was able to detect the sawtoothed grain beetle, \textit{O. surinamensis}, within barley and was able to determine the level of infestation using cross-correlation \cite{muller-blenkleMethodAcousticStorage2023}.

The Acoustic Emission Consulting AED 2010L was able to detect one to two insects per \si{\kilo\gram} of wheat grain with \SIrange{72}{100}{\percent} accuracy. The portable, battery-operated system uses a piezoelectric sensor mounted on a metal probe that is inserted into a grain sample. The system was able to predict the population densities of the tested live specimen including the rice weevil, \textit{S. oryzae}, the lesser grain borer, \textit{R. dominica}, the confused flour beetle \textit{T. confusum}, the sawtoothed grain beetle, \textit{O. surinamensis}, the rusty grain beetle, \textit{C. ferrugineus}, the khapra beetle, \textit{T. granarium}, and the cigarette beetle, \textit{L. serricorne}, through chi-squared tests and machine learning classifiers \cite{eliopoulosDetectionAdultBeetles2015,eliopoulosEstimationPopulationDensity2016}.

The Postharvest Insect Detection System (PDS) uses electret microphones centered on a plexiglass base onto which grain sample is placed. The system was tested with the rice weevil, \textit{S. oryzae}, and the red flour beetle, \textit{T. castaneum} \cite{mankinPerformanceLowCostAcoustic2020}. The Computerized Acoustical Larval Detection system was patented in 1988 and consists of a bowl capable of holding \SI{1}{\liter} of grain sample with a diaphram at the bottom that is coupled with a microphone through a conduit \cite{j.c.webb.c.a.litzkowComputerizedAcousticalLarval1988}. The Ultrasonic Insect Detector, also known as the Purdue INsect Feeding monitor, was patented in 1990 \cite{shadeDetectionHiddenInsect1990}. The active sensor was able to detect the cowpea beetle, \textit{Callosobruchus maculatus}, larva inside the cowpea seeds at the \SI{40}{\kilo\hertz} range \cite{shadeUltrasonicInsectDetector1989} and monitor the development of the rice weevil, \textit{Sitophilus o.} in maize kernels \cite{pittendrighMonitoringRiceWeevil1997}. The sensor is only effective at distances less than a few millimeters because of the attenuation of the sound waves at high frequencies when traveling through grain \cite{shadeUltrasonicInsectDetector1989}.

The Sensor Technology and Applied Research (STAR) Center at Stevens Institute of Technology developed and refined acoustic sensors and detection algorithms for identifying insects in stored products during transportation and import inspection. These systems were tested on live specimens, including the larger cabinet beetle, \textit{T. inclusum}, mealworm and darkling beetle, \textit{T. molitor}, corn earworm, \textit{H. zea}, \textit{Copitarsia} larvae, and khapra beetle, \textit{T. granarium}, across various substrates.

The initial platform featured a test container with ceramic piezoelectric sensors, delivering four audio channels to an external preamplifier and ADC. Housed within a sound-attenuated enclosure, the system achieved up to \SIrange{50}{80}{\deci\bel} attenuation above \SI{500}{\hertz}. Analysis showed insect signals were concentrated between \SIrange{400}{1000}{\hertz}. A processing chain separated recordings into noise and insect bands, extracted Hilbert envelopes, and used their difference as a detection metric, achieving up to 100\% accuracy for larvae and 70\% for adults, with no false positives in clean rice samples \cite{flynnAcousticMethodsInvasive2016}.

Later versions integrated electronics for portability and used a Pelican case lined with acoustic dampening material, providing \SI{60}{\deci\bel} attenuation above \SI{500}{\hertz}. The updated enclosure featured three channels with piezoelectric modules and a revised algorithm: bandpass filtering from \SIrange{500}{5000}{\hertz}, moving average smoothing, and normalization. Detection was indicated when the signal exceeded background noise by at least \SI{5}{\deci\bel} \cite{sutinVibroacousticMethodsInsect2017}.

\subsection{Acoustic Insect Datasets}

Several audio databases featuring insects have been published and are publicly available. The Animal Sound Archive \cite{registry-migration.gbif.orgAnimalSoundArchive2016}, InsectSet32 \cite{faissInsectSet32DatasetAutomatic2022}, InsectSet47 and InsectSet66 \cite{faissInsectSet47InsectSet66Expanded2023, faissAdaptiveRepresentationsSound2023}, InsectSound1000 \cite{brandingInsectSound1000InsectSound2024}, and Singing Insects of North America (SINA) \cite{SingingInsectsNorth2024} all include insect recordings; however, these datasets feature neither species commonly found infesting stored products nor audio recordings of insects within materials or stored product. BugBytes has 95 recordings of insects in different environments but only presents approximately 23 minutes of recordings with less than 15\% of that pertaining to stored product insects; only two of which recorded with piezoelectric disks \cite{mankinBugBytesSound2019}.

\newpage